\section{Thảo luận về scale tuyến tính và phi tuyến}

\subsection{Vì sao cần cả hai loại scale}

\textbf{Linear scale} phù hợp để đọc chênh lệch tuyệt đối về thời gian chạy (đơn vị ms).
Tuy nhiên khi một thuật toán lớn vượt trội, các đường còn lại dễ bị nén sát trục và khó quan sát.

\textbf{Log-$y$ scale} giúp:
\begin{itemize}[leftmargin=1.2cm]
    \item Giảm hiện tượng nén dữ liệu khi khoảng giá trị quá rộng.
    \item Làm rõ tốc độ tăng tương đối giữa các thuật toán.
    \item Dễ quan sát xu hướng trên nhiều bậc độ lớn.
\end{itemize}

\subsection{Nguyên tắc diễn giải đúng trên log scale}

Khi đọc Hình \ref{fig:time-logy}, cần lưu ý:
\begin{itemize}[leftmargin=1.2cm]
    \item Khoảng cách theo trục $y$ thể hiện tỷ lệ nhân, không phải chênh lệch cộng tuyến tính.
    \item Độ dốc trên log scale phản ánh mức tăng theo bậc độ lớn.
    \item Kết luận cuối cùng nên đối chiếu cả Hình \ref{fig:time-linear} và Hình \ref{fig:time-logy}.
\end{itemize}

\subsection{Đánh giá kỹ thuật trực quan đã sử dụng}

\begin{itemize}[leftmargin=1.2cm]
    \item \textbf{Faceting theo input type:} giảm nhiễu khi so sánh nhiều kiểu dữ liệu cùng lúc.
    \item \textbf{Màu + marker nhất quán:} hỗ trợ theo dõi cùng thuật toán xuyên suốt các hình.
    \item \textbf{Legend đặt ngoài vùng dữ liệu:} tránh che khuất đường hoặc cột.
    \item \textbf{Bản vector (PDF/SVG):} giữ độ sắc nét khi phóng to trong báo cáo và trình chiếu.
\end{itemize}
